\begin{theorem}[Horizontal movement in a common tree background]
    \label{thm:peps_gauged_horizontal_flux_move}
    \lean{
        TNLean.PEPS.IsGIsometric.exists_unitary_torusGaugedHorizontalFluxMove,
        TNLean.PEPS.torusGaugedHorizontalFluxAssignment,
        TNLean.PEPS.torusGaugedHorizontalFluxAssignment_one
    }
    \leanok
    On a regular torus with horizontal period at least four and vertical period
    at least three, let $R$ be the translated six-site horizontal two-plaquette
    block. Reconstruct both internal assignments from the same vertex labels
    $k$ and the single-bond or two-bond insertion of $g$. There is one unitary
    $W$ on the original spins in $R$, chosen before $k$, $g$, all boundary
    configurations and all common exterior and crossing operators, which sends
    the single insertion to the two-bond insertion. Its adjoint effects the
    reverse movement. The identity extension outside $R$ gives the corresponding
    identity of actual closed contractions. Both periodic seams are included.

    This is a finite-torus common-background specialization of the local movement
    construction in SCP10, Theorem~6.16, and its accessible-coordinate argument
    (local source lines 1765--1920 and 2271--2305). It does not identify different
    crossing presentations or establish the prescribed four-endpoint braid;
    see \cite{gap:rmp_peps_quantum_double_g_isometry}.
    \uses{thm:peps_regular_gauged_cycle_permutation}
\end{theorem}
\begin{proof}\leanok
    Use the actual translated induced graph, its fixed spanning path, and its two
    native vertical chords. The two-cycle permutation extends the first chord
    label to the second. The shared vertex labels give the same boundary
    reindexing on both sides. The original-spin operation therefore acts on all
    boundary columns, and the actual cut expansion gives the closed-state
    identity. Unitarity yields the reverse action.
\end{proof}
